\chapter*{解析几何每日一题}

\refer{《高等数学》同济版}

\exercise{
  求解以向量为元的线性方程组
  \formula{\left\{
    \addtolength{\jot}{-1em}
    \begin{aligned}
      & 5x-3y=a\\
      & 3x-2y=b
    \end{aligned}\right.
  }
  其中$\vv{a}=\qty(2,1,2),\ \vv{b}=\qty(-1,1,-2)$
  }

\teacher{成晖}

\postdate{2023年4月13日}

\result{
  \left\{
  \addtolength{\jot}{-1em}
  \begin{aligned}
    &x=\qty(7,-1,10)\\
    &y=\qty(11,-2,16)
  \end{aligned}\right.
}

\analysis{利用坐标对向量的线性运算，即对向量的各个坐标分别进行相应的数量运算}

\solution{
  如同解二元一次方程组，可得
  \formula{x=2a+3b,\ y=3a+5b}
  将$a,\ b$的坐标代入，得
  \formula{x=2\qty(2,1,2)-3\qty(-1,1,-2)=\qty(7,-1,10)}
  \formula{y=3\qty(2,1,2)-5\qty(-1,1,2)=\qty(11,-2,16)}
  本期的内容较为简单，在解析几何里，向量是一种最基本的名词，并且占据着不可或缺的地位，而向量的基本运算便更重要啦.
}